\section{Coherent distributions and entropy minima}
\label{ent:section}

\subsection{Normalization and entropy conventions}
\label{ent:definitions}

In this section, $(\mathcal O,\nu)$ denotes either the Slater orbit
with its invariant probability measure from Section~\ref{ext:section},
or a fixed-parity Gaussian orbit from Theorem~\ref{spin:convex-order}.
The representation dimension is respectively
$D=\binom np$ or $D=2^{n-1}$. The same definitions apply to the odd-spin
orbit in Corollary~\ref{spin:odd}. A unit coherent vector at $x\in\mathcal O$
is denoted by $\Omega_x$, and
$Q_\rho(x)=\langle\Omega_x,\rho\Omega_x\rangle$ is independent of its
phase.

The resolution identity gives
\begin{equation}
 \int_{\mathcal O}|\Omega_x\rangle\langle\Omega_x|\,d\nu(x)
 =\frac{\Id}{D},
 \qquad D\int_{\mathcal O}Q_\rho\,d\nu=1.
 \label{ent:resolution}
\end{equation}
For exterior powers this is \eqref{ext:resolution}. For the irreducible
spin representations, the averaged projector commutes with the group;
Schur's lemma and its trace one give the identity.

Define the unnormalized and normalized positive moments by
\begin{equation}
 M_\rho(q)=\int_{\mathcal O}Q_\rho^q\,d\nu,
 \qquad N_\rho(q)=D M_\rho(q),\qquad q>0.
 \label{ent:moments-definition}
\end{equation}
The normalized first moment is one. Since $Q_\rho\ge0$ has positive
integral, each $M_\rho(q)$ is strictly positive. All these moments are
finite because $0\le Q_\rho\le1$.

We use the resolution measure $D\,d\nu$ to define the
R\'enyi--Wehrl and Wehrl entropies:
\begin{align}
 S_q(\rho)&=\frac{\log N_\rho(q)}{1-q},
 &&q>0,\quad q\ne1,
 \label{ent:renyi-definition}\\
 S_W(\rho)&=-D\int_{\mathcal O}Q_\rho\log Q_\rho\,d\nu.
 \label{ent:wehrl-definition}
\end{align}
Logarithms are natural and $0\log0=0$. In particular, the latter
integrand is the continuous extension of $x\log x$ at zero.
The probability density relative to $\nu$ is $f_\rho=DQ_\rho$;
its differential entropy is
\begin{equation}
 -\int f_\rho\log f_\rho\,d\nu=S_W(\rho)-\log D.
 \label{ent:output-entropy}
\end{equation}

\subsection{The coherent Beta laws}

For $a,b>0$, the distribution $\operatorname{Beta}(a,b)$ on $(0,1)$
has density
\[
 \frac{\Gamma(a+b)}{\Gamma(a)\Gamma(b)}
 t^{a-1}(1-t)^{b-1},
\]
where $\Gamma(z)=\int_0^\infty s^{z-1}e^{-s}\,ds$ for $z>0$
is Euler's Gamma function. A $\operatorname{Gamma}(a,1)$ variable
has density $s^{a-1}e^{-s}/\Gamma(a)$ on $(0,\infty)$.
If $X$ and $Y$ are independent Gamma variables with shapes $a,b$
and common scale one, the change of variables
$(X,Y)=(rt,r(1-t))$ shows that $X/(X+Y)$ has distribution
$\operatorname{Beta}(a,b)$.

\begin{proposition}[Coherent measurement laws]
\label{ent:beta-products}
For a Slater input in $\mathcal H_{p,n}$ with $0<p<n$,
\begin{equation}
 Q_{\rho_W}\ \overset{\mathrm{law}}=
 \prod_{j=0}^{p-1} A_j,
 \qquad
 A_j\sim\operatorname{Beta}(p-j,n-p),
 \label{ent:exterior-beta}
\end{equation}
where the $A_j$ are independent. For $p=0$ or $p=n$, the coherent
measurement is identically one.

For a pure Gaussian input in either parity sector $\mathcal F_n^\epsilon$,
\begin{equation}
 Q_{\Omega}\ \overset{\mathrm{law}}=
 \prod_{k=1}^{n-1} B_k,
 \qquad B_k\sim\operatorname{Beta}(k,k),
 \label{ent:spin-beta}
\end{equation}
with independent factors. The product is one when $n=1$.
For the basic spin representation of $\Spin(2m+1)$, use
\eqref{ent:spin-beta} with $n=m+1$.
\end{proposition}
\begin{proof}
Let $Z$ be a standard complex Gaussian vector in $E=\C^n$, with
independent coordinates of density $\pi^{-1}e^{-|z|^2}$. Its direction
$u=Z/\|Z\|$ is uniform on the complex unit sphere. For a fixed
$p$-plane $W$, the variables
\[
 X=\|\pi_W Z\|^2,\qquad
 Y=\|(\Id-\pi_W)Z\|^2
\]
are independent Gamma variables of shapes $p$ and $n-p$.
Consequently the first conditional weight is
\begin{equation}
 a_{\rho_W}(u)=\|\pi_Wu\|^2=\frac{X}{X+Y}
 \sim\operatorname{Beta}(p,n-p).
 \label{ent:first-exterior-weight}
\end{equation}
For almost every $u$ this weight is positive. By
Lemma~\ref{ext:coherent-closure}, the normalized conditional input
is a Slater density in $\Lambda^{p-1}(u^\perp)$. Its inner Haar
Husimi distribution depends only on $(p-1,n-1)$, independently of
$u$. On the double-Haar probability space of
\eqref{ext:double-haar}, this conditional distribution being constant
means that the lower-dimensional readout is independent of the first
weight. Iterating \eqref{ext:conditioning} proves
\eqref{ent:exterior-beta}; at step $j$ the two Beta shapes are
$p-j$ and $(n-j)-(p-j)=n-p$.

For a Gaussian input, its Majorana covariance $\Gamma_\Omega$
satisfies $\Gamma_\Omega^2=-\Id$. A Haar-rotated first mode gives
a uniform oriented orthonormal two-frame $(u,v)$ in $\R^{2n}$.
Conditioned on $u$, the vector $v$ is uniform on the unit sphere in
$u^\perp$. The linear functional $v\mapsto u^{\mathsf T}\Gamma_\Omega v$
is represented there by the unit vector $\Gamma_\Omega^{\mathsf T}u$.
For $n\ge2$, its value $t$ therefore has density proportional to
\[
 (1-t^2)^{n-2},\qquad -1<t<1.
\]
The occupation identity \eqref{spin:occupation} gives
$a=(1+t)/2$, whose density is proportional to
$a^{n-2}(1-a)^{n-2}$. Thus
$a\sim\operatorname{Beta}(n-1,n-1)$.
The Gaussian conditional closure in Lemma~\ref{spin:gaussian-closure}
and the double-Haar identity \eqref{spin:double-haar} make the
remaining coherent readout independent of this weight, with the
law for $n-1$ modes. Induction proves \eqref{ent:spin-beta}.
The one-dimensional initial sector gives the empty product.
Finally, Corollary~\ref{spin:odd} identifies the entire odd-spin
coherent orbit and its probability measure with the indicated
half-spin orbit.
\end{proof}

The independence in Proposition~\ref{ent:beta-products} describes
coherent inputs. For a general density, its normalized conditional
state varies with the preceding measurement direction; the convex-order
proof retains that dependence.

\subsection{Positive moments and explicit minima}

The Beta integral gives
\begin{equation}
 \E A^q=
 \frac{\Gamma(a+q)\Gamma(a+b)}
 {\Gamma(a)\Gamma(a+b+q)},
 \qquad A\sim\operatorname{Beta}(a,b),\quad q>0.
 \label{ent:beta-moment}
\end{equation}
Consequently the coherent moments in the two series are
\begin{align}
 m_{p,n}(q)
 &:=\prod_{j=0}^{p-1}
 \frac{\Gamma(p-j+q)\Gamma(n-j)}
 {\Gamma(p-j)\Gamma(n-j+q)},
 \label{ent:exterior-moments}\\
 b_n(q)
 &:=\prod_{k=1}^{n-1}
 \frac{\Gamma(k+q)\Gamma(2k)}
 {\Gamma(k)\Gamma(2k+q)}.
 \label{ent:spin-moments}
\end{align}
Here $m_{0,n}(q)=1$ by the empty-product convention, while each
factor in $m_{n,n}(q)$ cancels to one. These formulas cover every
real $q>0$. In particular,
\begin{equation}
 m_{p,n}(1)=\binom np^{-1},\qquad
 b_n(1)=2^{-(n-1)},
 \label{ent:first-moments}
\end{equation}
as required by \eqref{ent:resolution}.

\begin{theorem}[Wehrl and R\'enyi--Wehrl minima]
\label{ent:minima}
Let $H_0=0$ and $H_r=\sum_{j=1}^r j^{-1}$ for $r\ge1$.
For every density on $\mathcal H_{p,n}$, $0\le p\le n$,
\begin{align}
 S_q(\rho)
 &\ge\frac{\log\!\left[\binom np\,m_{p,n}(q)\right]}{1-q},
 &&q>0,\quad q\ne1,
 \label{ent:exterior-renyi}\\
 S_W(\rho)
 &\ge s_{p,n}:=\sum_{j=0}^{p-1}(H_{n-j}-H_{p-j}).
 \label{ent:exterior-wehrl}
\end{align}
Equality in each inequality holds precisely for pure Slater densities.

For every density supported on the fixed parity sector
$\mathcal F_n^\epsilon$, $n\ge1$,
\begin{align}
 S_q(\rho)
 &\ge\frac{\log[2^{n-1}b_n(q)]}{1-q},
 &&q>0,\quad q\ne1,
 \label{ent:spin-renyi}\\
 S_W(\rho)
 &\ge s_n^{\mathrm{spin}}:=\sum_{k=1}^{n-1}(H_{2k}-H_k).
 \label{ent:spin-wehrl}
\end{align}
Equality holds precisely for pure Gaussian densities in that parity
sector. For the basic spin representation of $\Spin(2m+1)$, the
same formulas hold with dimension $2^m$, moment $b_{m+1}(q)$,
and constant $s_{m+1}^{\mathrm{spin}}$.
\end{theorem}
\begin{proof}
For $q>1$, the function $x\mapsto x^q$ is continuous and strictly
convex on $[0,1]$. Theorems~\ref{ext:convex-order} and
\ref{spin:convex-order} imply that its moment is maximized by a
pure coherent input. Multiplication by $D$, application of the
increasing logarithm, and division by $1-q<0$ give the entropy
lower bound. For $0<q<1$, use the continuous strictly convex
function $x\mapsto-x^q$. The moment inequality reverses, and
$1-q>0$, giving the same direction for the entropy inequality.
The positivity of the moments justifies every logarithm and its
equality criterion. Strict convexity identifies the equality cases
in both ranges.

For Wehrl entropy, apply the same theorems to the continuous strictly
convex function $x\mapsto x\log x$, with its value zero at the
origin, and multiply by $-D$. It remains to evaluate the coherent
integral. Differentiation of a moment at $q=1$ gives
\[
 \left.\frac{d}{dq}\int Q_\Omega^q\,d\nu\right|_{q=1}
 =\int Q_\Omega\log Q_\Omega\,d\nu.
\]
This differentiation is justified directly: for
$q\in[1/2,3/2]$ and $0<x\le1$,
$|x^q\log x|\le x^{1/2}|\log x|\le2/e$.
The derivative at $x=0$ is zero, so dominated convergence applies
on the probability space including all zeros of $Q_\Omega$.

Taking logarithmic derivatives in
\eqref{ent:exterior-moments} and \eqref{ent:spin-moments}, and
using $\Gamma(z+1)=z\Gamma(z)$, gives
\begin{align*}
 -D_{p,n}m_{p,n}'(1)
 &=\sum_{j=0}^{p-1}(H_{n-j}-H_{p-j}),\\
 -2^{n-1}b_n'(1)
 &=\sum_{k=1}^{n-1}(H_{2k}-H_k).
\end{align*}
For clarity, if $\psi_\Gamma=\Gamma'/\Gamma$, the recurrence
$\psi_\Gamma(r+1)=\psi_\Gamma(r)+1/r$ shows that each difference
$\psi_\Gamma(b+1)-\psi_\Gamma(a+1)$ at positive integers equals
$H_b-H_a$. Equation~\eqref{ent:first-moments} supplies the factor
$D$ in the preceding derivatives. This proves the displayed constants.
The odd-spin conclusions follow from Corollary~\ref{spin:odd}.
\end{proof}

\begin{example}[Small sectors]
\label{ent:examples}
For one particle in $n$ modes, the coherent Wehrl value is
$s_{1,n}=H_n-1$. For two particles in four modes,
\[
 s_{2,4}=(H_4-H_2)+(H_3-H_1)=\frac{17}{12}.
\]
The maximally mixed density on this six-dimensional sector has
$Q\equiv1/6$ and hence $S_W=\log6$, larger than the coherent
value. In the spin series the constants for $n=1,2,3$ are
$0$, $1/2$, and $13/12$. The empty and filled exterior sectors
have zero Wehrl and R\'enyi--Wehrl entropy.
\end{example}

All formulas above concern the basic exterior, half-spin, and odd-spin
representations specified in the corresponding theorems. The
representation parameter is not replaced by a higher Cartan power.
